\documentclass[12pt]{article}
\usepackage[utf8]{inputenc}
\usepackage[letterpaper,margin=1.0in]{geometry}
%% Some formatting stuff
\usepackage{authblk}
\usepackage{fancyhdr}
%\usepackage{lineno}
\usepackage{hyperref}
\usepackage{booktabs, array, longtable}
\usepackage{multirow}
\pagestyle{fancy}
\setlength{\headheight}{14.5pt} % Fix fancyhdr warning
% Custom citep command for biblatex
\newcommand{\citep}{\parencite}
\newcommand{\citet}{\textcite}
\fancyhead[R]{\textbf{pg\_gpu}}
% for figures
\usepackage{graphicx}
\usepackage{wrapfig}
\usepackage{float}
\usepackage{subcaption}
\usepackage{nameref}
\usepackage{amsmath}
\usepackage{amssymb} % \checkmark in tab:competitors
\usepackage{xspace}

\graphicspath{ {./figures/} }

\newcommand{\comment}[1]{{\color{blue} #1}}
\newcommand{\pggpu}{\texttt{pg\_gpu}\xspace}
% Software / package macros. Use these in prose so package names
% always typeset in \texttt and xspace handles the trailing space.
% Qualified forms (e.g. moments.LD, pandas.DataFrame) keep raw
% \texttt{...} since xspace breaks the dot path.
\newcommand{\vcftools}{\texttt{VCFtools}\xspace}
\newcommand{\scikitallel}{\texttt{scikit-allel}\xspace}
\newcommand{\pixy}{\texttt{pixy}\xspace}
\newcommand{\clam}{\texttt{clam}\xspace}
\newcommand{\plink}{\texttt{PLINK}\xspace}
\newcommand{\plinktwo}{\texttt{PLINK2}\xspace}
\newcommand{\angsd}{\texttt{ANGSD}\xspace}
\newcommand{\popgenome}{\texttt{PopGenome}\xspace}
\newcommand{\tskit}{\texttt{tskit}\xspace}
\newcommand{\moments}{\texttt{moments}\xspace}
\newcommand{\cupy}{\texttt{CuPy}\xspace}
\newcommand{\numpy}{\texttt{NumPy}\xspace}
\newcommand{\scipy}{\texttt{SciPy}\xspace}
\newcommand{\pandasname}{\texttt{pandas}\xspace}
\newcommand{\msprime}{\texttt{msprime}\xspace}
\newcommand{\zarr}{\texttt{zarr}\xspace}
\newcommand{\biotzarr}{\texttt{bio2zarr}\xspace}
\newcommand{\vcz}{\texttt{VCZ}\xspace}
\newcommand{\vcf}{\texttt{VCF}\xspace}
\newcommand{\pixi}{\texttt{pixi}\xspace}
\newcommand{\cublas}{\texttt{cuBLAS}\xspace}
\newcommand{\lostruct}{\texttt{lostruct}\xspace}

% for hyperlinks
\hypersetup{
    colorlinks=true,
    citecolor=black,
    urlcolor=cyan,
    linkcolor=blue
    }
\urlstyle{same}

% for highlighting text
\usepackage{xcolor}
\usepackage{soul}
\usepackage{tikz}

% code listings
\usepackage{listings}
\lstset{
    basicstyle=\ttfamily\small,
    breaklines=true,
    frame=single,
    language=Python
}

% bibliography
\usepackage[backend=biber,style=authoryear]{biblatex}
\addbibresource{refs.bib}

%\linenumbers
\renewcommand*{\bibfont}{\fontsize{10}{12}\selectfont}

\newcommand{\beginsupplement}{%
    \setcounter{table}{0}
    \renewcommand{\thetable}{S\arabic{table}}%
    \setcounter{figure}{0}
    \renewcommand{\thefigure}{S\arabic{figure}}%
    \renewcommand{\theHtable}{supp.\arabic{table}}%
    \renewcommand{\theHfigure}{supp.\arabic{figure}}%
}

\title{GPU accelerated population genetics statistics using \texttt{pg\_gpu}}
\author[1, 2]{Nathaniel S. Pope}
\author[1, 2]{Angel G. Rivera-Col\'{o}n}
\author[1, 2]{Ananya Kapoor}
\author[3]{Kevin Korfmann}
\author[4]{Murillo F. Rodrigues}
\author[1, 2, *]{Andrew D. Kern}
\affil[1]{\small{Institute of Ecology and Evolution, University of Oregon}}
\affil[2]{\small{Department of Biology, University of Oregon}}
\affil[3]{\small{Department of Biology, University of Pennsylvania}}
\affil[4]{\small{Oregon National Primate Research Center, Oregon Health \& Science University}}
\affil[*]{\small{Corresponding author: \href{mailto:adkern@uoregon.edu}{adkern@uoregon.edu}}}

\date{\small{\today{}}}

\begin{document}

\maketitle

\begin{abstract}
\noindent\textbf{Motivation:} Population genetics summary
statistics---diversity, divergence, linkage disequilibrium, selection
scans, and dimensionality reduction---are fundamental across human,
agricultural, and ecological genomics. As whole-genome sequencing
datasets have grown to hundreds of thousands of individuals, the
cost of computing these statistics on conventional CPU implementations has
become a major bottleneck: windowed scans of a single chromosome arm
can take hours to days, and computation of pairwise linkage-disequilibrium
statistics useful for demographic inference scales as $O(n^2)$ in
sample size, often exceeding wall-clock budgets entirely.

\noindent\textbf{Results:} We present \pggpu, a Python library
implementing a comprehensive catalog of population-genetics summary
statistics as fused CUDA kernels on NVIDIA GPUs. \pggpu covers eleven
categories spanning diversity and neutrality tests, divergence,
admixture, the site-frequency spectrum, linkage disequilibrium,
haplotype-based selection scans, dimensionality reduction (PCA,
randomized PCA, local PCA / \lostruct), distance distributions,
relatedness, resampling, and a generalized weighted-SFS framework for
custom $\theta$ estimators. On the full Ag1000G Phase 3 chromosome 3R
arm (2{,}940 haplotypes, 10.9 million variants) \pggpu agrees with
\scikitallel and \plinktwo to machine precision while delivering a
median $139\times$ and maximum $1{,}096\times$ speedup. For
the multi-population LD statistics used by \moments for demographic inference,
\pggpu is a drop-in replacement that yields a $\sim$1{,}750-fold
speedup over the native implementation. Whole chromosome arm scans, \lostruct screens, and calculation 
of LD statistics complete on a single NVIDIA A100 in
seconds to a few minutes.

\noindent\textbf{Availability and Implementation:} \pggpu is open
source, implemented in Python ($\geq$3.12) on top of \cupy, and
distributed via a \pixi environment. Source, documentation, and
tutorials are at \url{https://github.com/kr-colab/pg_gpu}.

\noindent\textbf{Contact:} \href{mailto:adkern@uoregon.edu}{adkern@uoregon.edu}

\noindent\textbf{Supplementary Information:} An extended description of
the design and implementation, the full statistic catalog with primary
references, an extended description of the API and features, numerical
accuracy comparisons against \scikitallel, \plinktwo, and \moments, a
worked \moments-LD demographic-inference example, and supplementary
figures (including a breakdown of the speed and accuracy of LD statistics parsing for moments) are
available at the journal website.
\end{abstract}

\section*{Introduction}

Population genetics relies on a rich toolkit of summary statistics to characterize
genetic variation within and between populations. Statistics such as nucleotide
diversity ($\pi$), Tajima's $D$, $F_{ST}$, and linkage disequilibrium (LD) are
foundational to studies of demographic history, natural selection, population
structure, and admixture. As whole-genome sequencing datasets have grown from
hundreds to hundreds of thousands of individuals---exemplified by projects such
as the Anopheles 1000 Genomes \citep{ag1000g2017} and UK Biobank
\citep{bycroft2018}---the computational cost of these analyses has become a
serious bottleneck.

A number of software tools have been developed to compute population genetics
statistics from genomic data. \plink \citep{purcell2007,chang2015} was the first 
high-performance tool to offer a broad range of genome-wide association as well as 
other simple population-based summaries. \vcftools \citep{danecek2011} was a
general-purpose tool that provided command-line access to diversity, $F_{ST}$,
and LD calculations from \vcf files, but was designed for the sample sizes of
the early 1000 Genomes era. \angsd \citep{korneliussen2014} is a
high-performance C/C++ tool, but it is oriented towards genotype likelihood 
workflows and consequently has limited support for population genetics summary 
statistics. \popgenome \citep{pfeifer2014} offered a
comprehensive R-based toolkit but has been removed from CRAN and is no longer
maintained. \scikitallel \citep{miles2024} brought population genetic analyses into 
the Python ecosystem with NumPy-based array operations, offering a broad repertoire of 
statistics including haplotype-based selection scans and $F$-statistics; however,
it is no longer under active development. \pixy \citep{korunes2021, bailey2025} 
addressed the important biases that missing data introduce in estimates in $\pi$, 
$d_{XY}$, $F_{ST}$, Tajima's $D$, and $\theta_W$, at the cost of computational 
performance. \clam is a very recent tool that focuses on faster computation of 
three statistics ($\pi$, $d_{XY}$, and $F_{ST}$) while accounting for missingness 
due to differences in callability across samples \citep{mirchandani2025}. Lastly,
\tskit \citep{jeffery2026} provides exceptionally efficient computation of a general 
class of statistics, which is enabled by a completely different data format, 
the tree sequence, rather than the traditional genotype matrix (for example, 
in \vcf or \plink formats). However, the tree sequences need to be first inferred 
from the genotype data itself, a process which may not be feasible for some datasets,
for example due to the lack of phased haplotypes.

Computation of population genetic statistics across all of these earlier tools is 
performed on the CPU, and statistics are typically computed sequentially along 
genomic windows. For large-scale projects involving windowed genome scans of tens 
of statistics across many populations, wall-clock times of hours to days are common. 
This limitation is particularly acute for two-locus statistics such as LD, where 
pairwise computations over variant sites are inherently $O(n^2)$ and form a bottleneck
 in inference pipelines such as \moments \citep{jouganous2017, ragsdale2019}.
Table~\ref{tab:competitors} summarizes the coverage and capabilities
of the most commonly used CPU tools alongside \pggpu (this work).

\begin{table}[ht]
\centering
\footnotesize
\caption{Coverage of population-genetics summary statistics across
established tools that operate on a genotype matrix (\vcf / \zarr) and
\pggpu. \checkmark, full support; $\sim$, partial or restricted
support; ---, no support or out of scope. Statistic categories follow
the broad families described in the text. \tskit \citep{jeffery2026}
is omitted because it operates on tree sequences rather than a
genotype matrix; \pggpu accepts both via \texttt{HaplotypeMatrix.from\_ts}
(rightmost column). \moments \citep{jouganous2017} is omitted because
its scope is downstream demographic inference rather than genome-scale
summary-statistic computation; \pggpu provides a drop-in replacement
for its LD-statistics parsing module (Supplementary Figure~\ref{fig:moments_ld}).}
\label{tab:competitors}
\setlength{\tabcolsep}{3pt}
\resizebox{\textwidth}{!}{%
\begin{tabular}{@{}lccccccccc@{}}
\toprule
\textbf{Tool}
 & \textbf{Backend}
 & \textbf{Diversity}
 & \textbf{Divergence}
 & \textbf{LD}
 & \textbf{Selection}
 & \textbf{PCA}
 & \textbf{Local PCA}
 & \textbf{GPU}
 & \textbf{Tree-seq.} \\
 & & \textbf{\& SFS} & \textbf{\& $F$-stats} & & \textbf{scans} & & \textbf{/ \lostruct} & & \textbf{input} \\
\midrule
\vcftools \citep{danecek2011}
 & C++         & $\sim$     & $\sim$     & $\sim$     & ---        & ---        & ---        & ---        & --- \\
\popgenome\textsuperscript{a} \citep{pfeifer2014}
 & R           & \checkmark & \checkmark & $\sim$     & $\sim$     & ---        & ---        & ---        & --- \\
\plinktwo \citep{chang2015}
 & C++         & $\sim$     & $\sim$     & $\sim$     & ---        & \checkmark & ---        & ---        & --- \\
\angsd \citep{korneliussen2014}
 & C++         & \checkmark & \checkmark & ---        & ---        & $\sim$     & ---        & ---        & --- \\
\scikitallel\textsuperscript{b} \citep{miles2024}
 & Python      & \checkmark & \checkmark & \checkmark & \checkmark & \checkmark & ---        & ---        & --- \\
\pixy \citep{korunes2021}
 & Python      & $\sim$     & $\sim$     & ---        & ---        & ---        & ---        & ---        & --- \\
\clam \citep{mirchandani2025}
 & Python      & $\sim$     & $\sim$     & ---        & ---        & ---        & ---        & ---        & --- \\
\pggpu (this work)
 & Python / CUDA & \checkmark & \checkmark & \checkmark & \checkmark & \checkmark & \checkmark & \checkmark & \checkmark \\
\bottomrule
\end{tabular}%
}
\par\smallskip
{\scriptsize\textsuperscript{a}Removed from CRAN; no longer
maintained. \textsuperscript{b}Active maintenance has ceased.}
\end{table}

Modern graphics processing units (GPUs) offer massive parallelism that is
well-suited to the data-parallel structure of population genetics computation.
Many statistics reduce to operations over matrices of haplotypes or genotypes
that can be decomposed into independent per-site or per-window calculations,
making them natural candidates for GPU acceleration. While GPUs have been
applied to specific problems in population genetics---notably forward-in-time
simulation \citep{gofish2017} and demographic inference via diffusion
approximations \citep{gutenkunst2021} or gradient-based optimization \citep{terhorst2025phlash} 
---no general-purpose GPU-accelerated library for computing population genetics summary 
statistics from empirical data has been developed.

Here, we present \pggpu, a Python library that leverages NVIDIA GPUs via \cupy
\citep{cupy2017} to compute a comprehensive catalog of population genetics
statistics. \pggpu implements over 70 functions spanning diversity, divergence,
the site frequency spectrum, linkage disequilibrium, haplotype-based selection
scans, admixture statistics, PCA, and relatedness---covering the large majority
of statistics in routine use. The library uses fused CUDA kernels --- kernels that compute
many related statistics in a single pass over the data --- and implements the
generalized $\theta$-estimation framework of \citet{achaz2009} to efficiently derive all
standard diversity and neutrality test statistics from a single site frequency
spectrum computation. \pggpu also provides a drop-in replacement for the LD-statistics parsing
module of \moments, achieving roughly $1{,}750$-fold speedups while
maintaining numerical accuracy at machine precision. Across windowed
genome scans, \pggpu delivers a median $139$-fold and up to
$1{,}096$-fold speedup relative to \scikitallel. \pggpu is
open source and freely available at
\url{https://github.com/kr-colab/pg_gpu}.

\section*{Design and Implementation}

% NSP 21-May-26: somewhere in this paragraph should mention the "streaming" design
% for large data?
\pggpu is a pure-Python library built on \cupy \citep{cupy2017}, with
around 50 custom CUDA kernels (\texttt{cupy.RawKernel}) and dense
linear algebra dispatched to \cublas; it depends only on \cupy,
\numpy \citep{harris2020}, \scipy \citep{virtanen2020}, \scikitallel
\citep{miles2024}, \msprime \citep{baumdicker2022}, and \zarr, and is
distributed as a \pixi environment so the full analysis stack installs
with one command. Two thin wrapper
classes---\texttt{HaplotypeMatrix} for phased haplotypes and
\texttt{GenotypeMatrix} for unphased diploid genotypes---carry data
through every analysis path; either can be built from a tree sequence
\citep{jeffery2026}, a \zarr store (\biotzarr/\vcz \citep{czech2024vcz}
or \scikitallel layouts), or an in-memory \numpy array, and moves
between CPU and GPU backings transparently.

The performance of \pggpu rests on a single observation: most
population-genetics statistics are per-site or per-pair reductions
that spend most of their CPU time reading the genotype matrix from
memory rather than doing arithmetic. \pggpu therefore uses
\emph{fused} kernels: one kernel reads each entry of the matrix once
and computes whichever statistics from a related family the user asked
for. A single entry point, \texttt{windowed\_analysis}, takes a list
of statistics and runs the fewest such kernels that together cover the
request. If the data are too large to fit in GPU memory at once, the
same kernels run in chunks whose size is picked automatically, so a
whole chromosome arm of $\sim$10 million variants fits on one A100
without any tuning by the user. Full details---the array layouts, 
the fused-kernel families, how missing data and accessibility masks 
are handled, the chunking rule, and the downstream interfaces---are 
described in the Supplementary Material (Extended Design and 
Implementation).

\section*{Features}

% NSP 21-May-26: anything relevant from the "streaming" design that is
% worth mentioning in passing here?
\pggpu provides a comprehensive catalog of population-genetics
statistics organized into eleven categories
(Table~\ref{tab:features}): classic estimators ($\pi$, $\theta_W$,
Tajima's $D$, Hudson and Weir--Cockerham $F_{ST}$, $d_{XY}$, the
unfolded and joint SFS, Patterson's $f_2$/$f_3$/$D$, iHS, nSL, Garud's
$H$) alongside several less commonly supported additions useful in
modern workflows---the Schrider et al.~\citep{schrider2018}
distance-based two-population statistics, the RAiSD
$\mu$ components~\citep{alachiotis2018}, the diplotype-frequency
spectrum, and a \texttt{FrequencySpectrum} interface to the
Achaz~\citep{achaz2009} weighted-SFS framework for custom $\theta$
estimators. Almost every statistic is also reachable through
\texttt{windowed\_analysis}, the entry point for genome
scans using the fused-kernel (Design and Implementation). A complete listing of all
implemented statistics with primary references, together with an
extended description of the API and resampling utilities, is provided
in the Supplementary Material (Table~\ref{tab:full_catalog}).

\begin{table}[ht]
\centering
\small
\caption{Summary of statistics implemented in \pggpu. The full catalog with
citations is provided in Supplementary Table~\ref{tab:full_catalog}.}
\label{tab:features}
\begin{tabular}{llp{5.5cm}}
\toprule
\textbf{Category} & \textbf{Count} & \textbf{Representative statistics} \\
\midrule
Diversity \& neutrality tests & 20+ & $\pi$, $\theta_W$, $\theta_H$, Tajima's $D$, Fay \& Wu's $H$, Zeng's $E$ \\
Divergence & 11 & $F_{ST}$ (Hudson; Weir--Cockerham; Nei), $d_{XY}$, $D_a$, PBS, $S_{nn}$, $G_{\min}$ \\
Linkage disequilibrium & 10+ & $r^2$, $D^2$, $D_z$, $\pi_2$, $Z_{nS}$, $\omega$, $\mu_{\mathrm{LD}}$, \texttt{moments.LD} parser \\
Selection scans & 9 & iHS, nSL, XP-EHH, XP-nSL, Garud's $H_{1/12/123}$, $H_{2}/H_{1}$, EHH decay \\
Site frequency spectrum & 6+ & Unfolded, folded, scaled, joint (2-pop) SFS \\
Admixture / $F$-statistics & 6+ & Patterson's $F_2$, $F_3$, $D$ (ABBA-BABA), block-jackknife wrappers \\
Dimensionality reduction & 5 & PCA, randomized PCA, PCoA, local PCA / \lostruct, corner detection \\
Relatedness / kinship & 2 & GRM (GCTA-equivalent), IBS (PLINK-equivalent) \\
Distance distribution & 4 & Hamming distance, variance, skewness, kurtosis \\
Resampling & 2 & Block jackknife, block bootstrap (incl.\ ratio-of-sums) \\
Generalized $\theta$ framework & -- & \texttt{FrequencySpectrum} class: custom weights, paired neutrality tests, SFS projection \\
\bottomrule
\end{tabular}
\end{table}

\section*{Results}

\subsection*{Benchmark dataset}

For benchmarking, we use Phase 3 of the \emph{Anopheles
gambiae} 1000 Genomes Project (Ag1000G) \citep{ag1000g2017}, a
deep-coverage whole-genome resequencing panel of 1{,}470 phased
mosquitoes (2{,}940 haplotypes) sampled from 22 populations across
sub-Saharan Africa. Chromosome arm 3R alone carries 10.9 million
phased, biallelic SNPs after filtering and accessibility masking; the
53\,Mb arm contains multiple known segregating inversions and
well-characterized insecticide-resistance sweeps, and is among the
most variable autosomal regions sequenced in any eukaryote. The
combination of extreme levels of segregating variation, non-trivial sample size, 
natural missingness, and rich population structure makes it a stringent and
biologically-meaningful stress-test for the fused CUDA kernels. Where a ground truth is
needed---for missing-data bias, runtime scaling, and
demographic-inference cross-checks against \moments---we pair
Ag1000G with \msprime \citep{baumdicker2022} simulations so each
result is supported by both empirical and simulation evidence.

\subsection*{Accuracy, performance, and scaling}

We first verified that \pggpu reproduces results from established CPU
implementations: every statistic with a counterpart in \scikitallel,
\plinktwo, or \moments matches the reference to within floating-point
error, with per-statistic comparisons reported in the Supplementary
Material (Table~\ref{tab:accuracy}, Supplementary
Figures~\ref{fig:plink_pca} and~\ref{fig:moments_ld_decay}).

Across the 23 Ag1000G 3R statistics with a directly comparable
\scikitallel implementation (2{,}940 haplotypes, 10.9 million
variants), the median \pggpu speedup is $139\times$ and the maximum
$1{,}096\times$ (Weir--Cockerham $F_{ST}$; Figure~\ref{fig:speedups},
Supplementary Figure~\ref{fig:walltimes}): divergence and pairwise-LD
kernels gain the most ($600$--$1{,}100\times$), single-population
diversity and SFS statistics $130$--$380\times$, and haplotype-based
selection scans $9$--$56\times$, the last bounded by per-haplotype
work that \scikitallel's C extensions already vectorize. On a
50\,kb-windowed slice of the same data, \pggpu and \scikitallel
estimates of $\pi$, an $R_0$-style heterozygosity ratio, Tajima's $D$,
and LD decay overlay to within line width while the windowed scan and
LD-decay computation run $5.8\times$ and $323\times$ faster
(Supplementary Figure~\ref{fig:overview}); a controlled cross-check on simulated
two-population \msprime data reproduces the same speedup pattern
against both \scikitallel and \plinktwo
(Supplementary Figure~\ref{fig:simbench}).

\begin{figure}[ht]
\centering
\includegraphics[width=0.85\textwidth]{benchmark_speedups.pdf}
\caption{\pggpu speedup over \scikitallel on Ag1000G Phase 3 chromosome
3R (2{,}940 haplotypes, 10.9M variants). The all-pairwise Rogers--Huff
LD comparison is run on the first $10{,}000$ SNPs (200 haplotypes) to
keep the \scikitallel reference tractable; all other statistics are
on the full chromosome arm. Bars show wall-clock speedup on a single
NVIDIA A100 GPU; orange bars indicate statistics where \scikitallel's
existing optimized implementation limits the achievable gain
(typically $<10\times$).}
\label{fig:speedups}
\end{figure}

A separate bottleneck is the tallying of multi-population two-locus haplotypes into LD statistics 
used for demographic inference \citep{ragsdale2019}. On a four-population coalescent
simulation (105 LD statistics, 10 recombination-distance bins, 10
replicates of 1\,Mb) \pggpu agrees with \moments-LD to a maximum
relative error of $1.8\times 10^{-11}$ while cutting per-replicate
parsing time from $256$\,s to $0.2$\,s---a $1{,}756 \pm 410\times$
speedup (Supplementary Figure~\ref{fig:moments_ld})---making
whole-genome LD-based inference tractable on standard hardware; a
per-population-count breakdown is given in Supplementary
Figure~\ref{fig:moments_ld_breakdown}.

Runtime scaling was characterized along two axes with \msprime
simulations \citep{baumdicker2022}: sample size (100 to 100{,}000
haplotypes at a fixed 100{,}000 variants) and variant count ($\sim
10^3$ to $\sim 10^6$ segregating sites at a fixed 200 haplotypes, with
realistic LD; Supplementary Figure~\ref{fig:scaling}). Per-site
statistics scale linearly on both axes and finish in well under a
second across the tested range; two-locus statistics show the expected
% NSP 21-May-26: why is Garud's H mentioned here as a two-locus statistic?
$O(n^2)$ growth with sample size for Garud's $H$ and a sub-quadratic
profile for $Z_{nS}$ (a sampled-pairs estimator). At the largest
configuration tested (100{,}000 haplotypes $\times$ 100{,}000 variants)
every single-pass statistic completes in under ten seconds on one A100.

\subsection*{Benchmarks and Applications}

We illustrate two end-to-end uses of \pggpu on data and methodology
where the GPU implementation enables an analysis that is slow or
intractable on CPU. A third use---supplying \pggpu-computed LD
statistics to the \moments demographic-inference workflow as a
single-import drop-in replacement, with downstream parameter estimates
indistinguishable from an analysis using the native parser---is described in the
Supplementary Material (Supplementary Figure~\ref{fig:moments_inference}).

\paragraph{Benchmarking a whole-arm Ag1000G genome scan.}
Partitioning the 1{,}470 Ag1000G samples into West and East African
groups (200 diploids / 400 phased haplotypes each) and running a full population-genetic
scan of chromosome 3R---loading 10.9 million phased variants from
\zarr, attaching the canonical accessibility mask, transferring to
GPU, and computing 19 windowed statistics in 100\,kb / 10\,kb
windows plus per-population and joint SFS---takes 168 seconds total
wall time, of which only 30 seconds is GPU compute
(Figure~\ref{fig:scan}). The same workflow expressed in
\scikitallel runs in roughly 75 minutes (windowed statistics alone),
so a single A100 reduces the runtime from over an hour
to seconds while recovering the expected genome-wide
signatures: $\pi$ heterogeneity, negative Tajima's $D$ consistent
with non-equilibrium demography and selection, and elevated Hudson $F_{ST}$ between subspecies.
Every per-window panel comes from a single call to \pggpu's \texttt{windowed\_analysis}
function, which will calculate multiple statistics (e.g., this plot) 
from the loaded haplotype matrix in a single forward pass.

\begin{figure}[ht]
\centering
\includegraphics[width=\textwidth]{ag1000g_genome_scan.png}
\caption{Multi-statistic scan of Ag1000G chromosome 3R using \pggpu.
Left column, eight per-window panels in 100\,kb windows with a 10\,kb
step ($H_{12}$ uses 1{,}000-SNP windows on its own grid):
nucleotide diversity $\pi$, Watterson's $\theta_W$, Tajima's $D$,
Fay \& Wu's $H^*$, and Zeng's $E$ (all West Africa); Hudson $F_{ST}$
and $d_{XY}$ between West and East Africa; and Garud's $H_{12}$
(West Africa). The red band on every left-column panel marks the
\textit{Gste} locus (3R:28.48--28.60\,Mb), a well-characterised insecticide-resistance
locus that the scan recovers as coordinated drops in $\pi$ and
$\theta_W$ with elevated $H_{12}$. Right column: West African
unfolded SFS, joint SFS heatmap (West vs.\ East), and per-step
compute time. Total compute time is 30\,s on a single A100 GPU;
total wall time including I/O is 168\,s.}
\label{fig:scan}
\end{figure}

\paragraph{Whole-genome local PCA.}
The local-PCA / \lostruct \citep{li2019} analysis workflow---comprised of
per-window local PCA, MDS using Frobenius distance between window covariances,
and corner detection---identifies genomic intervals whose structure deviates
from the genome-wide pattern. As a sanity check with a controlled signal we
applied it to a coalescent simulation \citep{baumdicker2022} containing a
single hard sweep (Supplementary Figure~\ref{fig:lostruct}): the embedding cleanly separates
neutral, linked, and sweep windows and the leading MDS axis tracks the
localized peak of Garud's $H_{12}$ at the sweep focal site.
Applied to the West African subset of Ag1000G 3R (100 diploids / 200 phased
haplotypes, 10.9M variants, $\sim$10{,}900 non-overlapping $1{,}000$-SNP
windows; Figure~\ref{fig:lostruct_ag1000g}), three distinct corner clusters are
recovered, and the corner-1 region near $\sim$28\,Mb overlaps with a localized
peak in the companion Garud's $H_{12}$ track (the \textit{Cyp6}/\textit{Gste}
insecticide-resistance locus on 3R), a known sweep target in West African
\emph{An.\ gambiae}.

\begin{figure}[ht]
\centering
\includegraphics[width=0.95\textwidth]{ag1000g_lostruct.pdf}
\caption{\lostruct on the West African subset of Ag1000G 3R (100
diploids / 200 phased haplotypes, 10.9M variants, $\sim$10{,}900 non-overlapping
$1{,}000$-SNP windows). Left: MDS embedding of windowed PCA
distances; windows in each of three detected corner clusters are
circled. Right: \lostruct MDS axis 1 along the chromosome (top)
and companion Garud's $H_{12}$ track on the same windows (bottom).
}
\label{fig:lostruct_ag1000g}
\end{figure}

\paragraph{Streaming a biobank-scale genome scan.}
At biobank sample sizes (e.g. $\ge 10^5$ individuals) a single chromosome will not fit on a GPU
all at once. \pggpu handles this by reading the chromosome one
chunk at a time; the same windowed-statistic, SFS, LD, and
pairwise-relatedness functions used on genotype matrices that can fit in memory work
without change on the chunked input. We demonstrate this on
chromosome 15 of \texttt{stdpopsim}'s \texttt{OutOfAfrica\_2T12}
model simulated at 50{,}000 diploid individuals per population
% NSP 21-May-26: wow! is the 2.3 TB a compressed or flat text VCF? If the former,
% we should clarify
(200{,}000 haplotypes, 11.6 million variants, $\sim 2.3$ TB on disk as a VCF).
The full scan shown in Supplementary
Figure~\ref{fig:biobank_scan}---windowed diversity and divergence
across all 200{,}000 haplotypes at three window sizes, marginal
SFS per population, a joint SFS projected from the full panel
to a smaller grid via per-variant hypergeometric sampling,
plus Garud's $H$ and an LD decay curve on smaller per-population
subsamples---finishes in about 16 minutes on a single A100 80\,GB. 
GPU memory use stays around 12\,GB regardless of how long the chromosome is.

\section*{Discussion}

\pggpu fills a long-standing gap in the population-genetics software
ecosystem. Existing tools optimize specific axes---\vcftools and
\plink for command-line access to focused sets of statistics,
\scikitallel for a Python-native research workflow, \pixy for
unbiased estimation under missingness, and \tskit for tree-sequence inputs,
but none brings the full catalog of routine summary
statistics onto modern GPU hardware. 
By unifying eleven categories of
statistics under a single Python API on top of \cupy, \pggpu turns
% NSP 21-May-26: does "batch" mean chunked-and-CPU-parallelized here? Or would "serial" be a better word
analyses that were previously hours-long batch jobs (e.g. full chromosome arm scans using
\scikitallel) or computationally infeasible
(chromosome-scale LD calculation using all SNP pairs) into runs of seconds to minutes on a single GPU.

% NSP 8-May-26: this point is great, and might be worth hinting at more heavily earlier
Three design choices that emerged from this work may be useful to
other authors of GPU-accelerated scientific software. First, fused
kernels that emit any subset of a related family of statistics from
registers in a single launch substantially outperform sequential
per-statistic GPU calls, because the dominant cost on this class of
workload is reading the haplotype matrix off-RAM rather than the
arithmetic itself. 
% NSP 8-May-26: sentence below seems like a given when there's a finite buffer---
% is there more to it?
Second, memory-aware chunking that
routes large workloads through a
chunked variant of the same kernels which is 
essential to handle large-scale data without user
tuning. 
% NSP 21-May-26: ^^^ this seems like a good place to mention any insights from the new "streaming" structure?
Third, processing windows in tiles---so that working memory
scales with the tile size ($O(\text{tile\_size}\times n_{\text{hap}}^2)$)
rather than with the number of windows ($O(n_{\text{windows}}\times n_{\text{hap}}^2)$)---lets
the same code path scale from a single chromosome to a whole genome.

Several limitations are worth stating. \pggpu requires an NVIDIA GPU
because \cupy targets CUDA; ports to AMD ROCm or Apple Metal would
require non-trivial kernel rewrites. Haplotype-pattern statistics
that scale as $O(n_{\text{hap}}^2)$ in sample size remain bound by
GPU memory at large scales (we tested up to 100{,}000 haplotypes);
genuine biobank-scale sample-sizes will need additional chunked-pair
kernels for the haplotype-pair statistics. The randomized SVD
used for the local PCA is an approximation to the
decomposition of the full covariance matrix;
although a dense eigendecomposition is available when exact \lostruct outputs
are required.
Finally, \pggpu currently supports phased haplotypes and unphased
diploid genotypes but does not yet handle low-coverage genotype
likelihoods natively (cf.\ \angsd); extending the framework to
genotype likelihoods is straightforward in principle and is
part of planned future work.

Looking forward, we see \pggpu as a foundation for two further
directions: integration with simulation-based machine-learning
frameworks for demographic and selection inference, where rapid
summary-statistic computation is the bottleneck during training;
and as a reference implementation of the Achaz \citep{achaz2009}
weighted-SFS framework, lowering the barrier for population
geneticists to design and validate custom test statistics. We expect
both will broaden the set of analyses that fit comfortably within
standard whole-genome research workflows.

\section*{Availability}

\pggpu is freely available at \url{https://github.com/kr-colab/pg_gpu}.

\section*{Acknowledgments}

We thank members of the Kern-Ralph co-lab for feedback on the library and the manuscript.

\section*{Funding}

This work was supported by NIH grants R01HG010774, R01HG012473, and R35GM148253.


\section*{Generative AI use}

Anthropic's Claude (Opus 4.6) was used to assist in drafting and
optimizing the fused \texttt{cupy.RawKernel} CUDA implementations
underlying \pggpu. All model-generated kernel code was reviewed,
profiled, and validated by the authors against appropriate
reference implementations (Supplementary Table~\ref{tab:accuracy})
before inclusion in the final codebase.



\printbibliography

\clearpage
\beginsupplement

\section*{Supplementary Material}

\subsection*{Numerical accuracy}

We validated \pggpu against \scikitallel using a 4\,Mb region of Ag1000G
chromosome 3L (3L:1{,}000{,}000--5{,}000{,}000; 2{,}940 haplotypes after
removing missing-data sites). Statistics whose estimators reduce to
exact integer counts---site-frequency spectra, two-locus
haplotype-frequency tallies (Garud's $H$), and Patterson's
$f_2$---agree exactly. Statistics that involve harmonic sums or
unbiased corrections ($\theta_W$, Tajima's $D$) agree to a relative
error of at most $5\times 10^{-4}$ (Table~\ref{tab:accuracy}). A
directly comparable Hudson $F_{ST}$ computation against \plinktwo on
simulated \msprime data agrees to $1.7\times 10^{-8}$, and the
corresponding PCA scores agree to $|r| > 0.88$ through PC5
(Figure~\ref{fig:plink_pca}). The two-population LD-decay curves
estimated by \pggpu and \moments on identical replicates are likewise
numerically indistinguishable across all nine LD statistics shown
(Figure~\ref{fig:moments_ld_decay}).

\begin{table}[ht]
\centering
\small
\caption{Agreement between \pggpu and \scikitallel on Ag1000G 3L
(2{,}940 haplotypes, 436{,}555 biallelic sites). Scalar statistics report
the relative error $|x_{\text{pg}} - x_{\text{al}}| / |x_{\text{al}}|$.
Vector statistics (SFS, windowed scans) report the Pearson correlation
across $n$ entries; for windowed statistics the median per-window
relative error is also given.}
\label{tab:accuracy}
\begin{tabular}{llrrr}
\toprule
\textbf{Category} & \textbf{Statistic} & \textbf{\pggpu} & \textbf{\scikitallel} & \textbf{Rel. error} \\
\midrule
\multirow{3}{*}{Diversity}
 & $\pi$ & $1.8660 \times 10^{-3}$ & $1.8660 \times 10^{-3}$ & $2.5 \times 10^{-7}$ \\
 & $\theta_W$ & $3.2257 \times 10^{-3}$ & $3.2250 \times 10^{-3}$ & $2.1 \times 10^{-4}$ \\
 & Tajima's $D$ & $-1.3642$ & $-1.3635$ & $5.1 \times 10^{-4}$ \\
\midrule
\multirow{3}{*}{Divergence}
 & Hudson $F_{ST}$ & $2.9060 \times 10^{-3}$ & $2.9060 \times 10^{-3}$ & $2.8 \times 10^{-15}$ \\
 & Weir--Cockerham $F_{ST}$ & $2.9262 \times 10^{-3}$ & $2.9262 \times 10^{-3}$ & $1.0 \times 10^{-15}$ \\
 & $d_{XY}$ & $1.9002 \times 10^{-3}$ & $1.9002 \times 10^{-3}$ & $1.1 \times 10^{-16}$ \\
\midrule
\multirow{2}{*}{Selection}
 & Garud's $H_1$ & $5.000 \times 10^{-3}$ & $5.000 \times 10^{-3}$ & $0$ \\
 & Garud's $H_{12}$ & $5.050 \times 10^{-3}$ & $5.050 \times 10^{-3}$ & $0$ \\
\midrule
Admixture & Patterson's $F_2$ (mean) & $2.278 \times 10^{-4}$ & $2.278 \times 10^{-4}$ & $0$ \\
\midrule
\multirow{2}{*}{SFS (Pearson $r$)}
 & Unfolded SFS, $n{=}201$ & --- & --- & $r = 1.0000$ \\
 & Joint SFS, $n{=}5{,}174$ & --- & --- & $r = 1.0000$ \\
\midrule
\multirow{3}{*}{Windowed (50\,kb, $n{=}80$)}
 & $\pi$ & --- & --- & $4.5 \times 10^{-15}$ \\
 & $\theta_W$ & --- & --- & $1.6 \times 10^{-4}$ \\
 & Hudson $F_{ST}$ & --- & --- & $1.9 \times 10^{-15}$ \\
\bottomrule
\end{tabular}
\end{table}

\clearpage

\begin{figure}[ht]
\centering
\includegraphics[width=0.85\textwidth]{benchmark_walltimes.pdf}
\caption{Absolute wall-clock time on Ag1000G chromosome 3R for the
23 statistics with both \pggpu and \scikitallel implementations.
Companion to Figure~\ref{fig:speedups}.}
\label{fig:walltimes}
\end{figure}

\begin{figure}[ht]
\centering
\includegraphics[width=0.9\textwidth]{benchmark_simulated.pdf}
\caption{Performance comparison on simulated \msprime data
(two-population split, $N_e = 10^4$, 10\,Mb sequence,
500 diploid individuals per population) against \scikitallel and
\plinktwo on shared statistics. The simulated benchmark provides a
controlled cross-check of the Ag1000G results in
Figure~\ref{fig:speedups}.}
\label{fig:simbench}
\end{figure}

\begin{figure}[ht]
\centering
\includegraphics[width=0.95\textwidth]{scikit_allel_comparison.pdf}
\caption{Equivalence and speedup of \pggpu vs.\ \scikitallel on a
single chromosome of phased Ag1000G data. Top three panels:
nucleotide diversity ($\pi$), $R_0$, and Tajima's $D$ in 50\,kb
windows along the chromosome, computed by both tools---curves
overlay to within line width. Lower panel: LD-decay (mean $r^2$ vs
pairwise distance) showing identical estimates from \scikitallel
and the two \pggpu LD estimators (Rogers--Huff and unbiased
$\sigma_d^2$). Horizontal bars at the bottom: wall-clock times for
the windowed scan ($5.8\times$ faster) and LD-decay computation
($323\times$ faster).}
\label{fig:overview}
\end{figure}

\begin{figure}[ht]
\centering
\includegraphics[width=\textwidth]{scaling_combined.pdf}
\caption{Runtime scaling with sample size (left, fixed at 100{,}000
variants) and variant count (right, fixed at 200 haplotypes). Both
axes are log-scaled. \texttt{windowed\_3} is a fused windowed kernel
that computes $\pi$, $\theta_W$, and Tajima's $D$ in a single pass.}
\label{fig:scaling}
\end{figure}

\begin{figure}[ht]
\centering
\includegraphics[width=0.85\textwidth]{accuracy_vs_plink_pca.pdf}
\caption{Per-axis correspondence of \pggpu PCA scores against \plinktwo
on simulated data (800 haplotypes, 19{,}470 biallelic variants,
two-population split). PC1 agrees to $|r| = 0.9997$; subsequent axes
show progressively reduced absolute correlation that reflects the
expected sensitivity of low-eigenvalue eigenvectors to floating-point
roundoff, with $|r| > 0.88$ through PC5.}
\label{fig:plink_pca}
\end{figure}

\begin{figure}[ht]
\centering
\includegraphics[width=0.85\textwidth]{missing_data_bias.pdf}
\caption{Bias of \pggpu's nucleotide-diversity estimator under
randomly induced missing genotypes (200 \msprime replicates per
missing-data rate). The estimator is unbiased across missing rates
from 0 to 60\%, mirroring the per-site valid-sample-count handling
implemented for all single-population diversity estimators.}
\label{fig:missing_bias}
\end{figure}

\begin{figure}[ht]
\centering
\includegraphics[width=0.85\textwidth]{moments_ld_validation.pdf}
\caption{\pggpu The accuracy and speed of parsing LD statistics, against
\moments. Left: raw two-locus LD-statistic sums ($D^2$, $Dz$,
$\pi_2$, $H$) for a four-population split simulation (10 replicates;
105 statistics). Right: per-replicate parsing time on
\moments CPU vs.\ \pggpu GPU ($1{,}756 \pm 410\times$ speedup).}
\label{fig:moments_ld}
\end{figure}

\begin{figure}[ht]
\centering
\begin{minipage}{\textwidth}
\centering
\includegraphics[width=0.95\textwidth]{moments_ld_2pop.png}\\
\includegraphics[width=0.95\textwidth]{moments_ld_3pop.png}\\
\includegraphics[width=0.95\textwidth]{moments_ld_4pop.png}
\end{minipage}
\caption{Numerical validation of 2-, 3-, and 4-population LD
statistics on a 500\,kb bp-binned stress test ($\mu=10^{-7}$, no
recombination map). \pggpu agrees with \moments to machine
precision; reported speedups ($193\times$, $288\times$, $246\times$)
are lower than Figure~\ref{fig:moments_ld} because that benchmark
uses the slower r-binned \moments path on $20\times$ more
sequence.}
\label{fig:moments_ld_breakdown}
\end{figure}

\begin{figure}[ht]
\centering
\includegraphics[width=0.85\textwidth]{moments_ld_decay_comparison.pdf}
\caption{Two-population LD decay curves estimated by \pggpu (this work)
and \moments on identical replicate simulations. The two
parsers produce numerically identical outputs across all 9 LD
statistics shown, confirming \pggpu can serve as a transparent
drop-in replacement for the native \texttt{moments.LD} parsing module.}
\label{fig:moments_ld_decay}
\end{figure}

\begin{figure}[ht]
\centering
\includegraphics[width=0.95\textwidth]{local_pca.png}
\caption{Local PCA (\lostruct) on a coalescent simulation containing a
hard sweep at 5\,Mb. Left: MDS embedding of windowed PCA distances,
colored by k-means regime (neutral, linked, sweep). Right: MDS axis
1 (top) and Garud's $H_{12}$ (bottom) along the chromosome.}
\label{fig:lostruct}
\end{figure}

\begin{figure}[ht]
\centering
\includegraphics[width=\textwidth]{genome_scan_ooa.pdf}
\caption{Biobank-scale genome scan of chromosome 15 from
\texttt{stdpopsim}'s \texttt{OutOfAfrica\_2T12} model, simulated
at 50{,}000 diploids per population (200{,}000 haplotypes,
11.6 million variants). Left column shows the per-window scan
across all 100{,}000 haplotypes per population: nucleotide
diversity $\pi$, Watterson's $\theta_W$, Tajima's $D$, and
normalized Fay \& Wu's $H^*$ in 100\,kb windows; Hudson $F_{ST}$
and $d_{XY}$ between populations; and Garud's $H_{12}$ in 10\,kb
windows on a 1{,}000-haplotype subsample per population (the
Garud kernel caps near 1{,}024 haplotypes). The red band marks
the 1\,Mb sub-region used for the pairwise-$r^2$ heatmap. Right
column, top to bottom: joint SFS projected from the full
100{,}000-haplotype panel per population to a 200-by-200 display
grid via per-variant hypergeometric sampling (the full
100{,}000-by-100{,}000 histogram is both unviewable and larger
than the GPU, but every variant from every haplotype contributes
to the projected version); mean $r^2$ vs.\ distance between
SNPs on common variants (minor-allele frequency $\geq 0.15$,
5{,}000-haplotype subsample per population) pooled across 16
probe regions tiling the chromosome; and the pairwise-$r^2$
heatmap of the 1\,Mb sub-region (5{,}000-haplotype subsample,
single population) with a zoom inset on the densest LD block.
The complete scan finishes in about 16 minutes on a single A100
80\,GB; GPU memory use stays around 12\,GB regardless of
chromosome length.}
\label{fig:biobank_scan}
\end{figure}

\clearpage

\subsection*{Demographic inference with \pggpu LD statistics}

Because \pggpu's \texttt{compute\_ld\_statistics} returns LD-statistic
sums in \moments' canonical schema (Design and Implementation), it can
replace the native \texttt{moments.LD} parser in an existing inference
pipeline by changing a single import. We verified this end to end on a
three-population isolation-with-migration model: fitting the model with
\pggpu-calculated LD statistics inside the standard \moments workflow
recovers parameter estimates indistinguishable from those obtained with
the native parser (Figure~\ref{fig:moments_inference}), while computing
the input statistics at the speedups reported in
Figure~\ref{fig:moments_ld}.

\begin{figure}[ht]
\centering
\begin{minipage}{\textwidth}
\centering
\includegraphics[width=0.2\textwidth]{moments_3pop_true_model.png}
\subcaption{Three-population isolation model used for simulation.}
\end{minipage}

\begin{minipage}{\textwidth}
\centering
\includegraphics[width=0.45\textwidth,keepaspectratio]{moments_3pop_fitted_vs_observed.png}
\subcaption{LD decay curves under the inferred model (\pggpu-calculated,
blue; \moments-parsed, green; truth, dashed orange) and
empirical bin estimates (black points). Inference using the two
parsers is indistinguishable.}
\end{minipage}
\caption{Demographic inference using \pggpu as the
\texttt{moments.LD} parser. The drop-in replacement preserves
parameter estimates while yielding the speedups in
Figure~\ref{fig:moments_ld}.}
\label{fig:moments_inference}
\end{figure}

\clearpage

\subsection*{Extended Design and Implementation}

\subsubsection*{Core data abstractions}

Two thin wrapper classes carry data through every analysis path. The
\texttt{HaplotypeMatrix} stores phased haplotypes as an
$n_{\mathrm{hap}}\times n_{\mathrm{var}}$ 8-bit integer array (with $-1$
denoting a missing genotype), variant positions, and optional named
sample sets. The companion \texttt{GenotypeMatrix} stores unphased
diploid genotypes as an $n_{\mathrm{samp}}\times n_{\mathrm{var}}$ 
8-bit unsigned-integer array---a layout chosen so that kernels whose
arithmetic factorizes over the genotype tabulation $(n_{AA}, n_{Aa},
n_{aa})$ read contiguous memory. Both classes retain a single canonical
buffer per object and route between CPU (\numpy) and GPU (\cupy)
backings transparently via a \texttt{device} property, so the same call
works whether the caller has a tree sequence \citep{jeffery2026}
(\texttt{HaplotypeMatrix.from\_ts}), a \zarr store (\texttt{from\_zarr},
supporting both \biotzarr/\vcz \citep{czech2024vcz} and \scikitallel
\citep{miles2024} layouts), or an in-memory \numpy array
\citep{harris2020}.

% NSP 21-May-26: seems like we need some details here regarding the "streaming" variants?

\subsubsection*{Fused kernels}

Most population-genetics statistics decompose into per-site or per-pair
reductions whose CPU implementations spend most of their time moving
data rather than computing: each pass over the haplotype matrix reads
tens of gigabytes of memory to emit a handful of summary scalars.
\pggpu takes advantage of this by writing \emph{fused} kernels that
read each input element once and emit \emph{any subset} of a related
family of statistics in a single pass. The entry point \texttt{windowed\_analysis}
accepts a list of statistic names and dispatches the request to the
smallest covering set of fused kernels: a single-population kernel
produces any combination of $\pi$, $\theta_W$, $\theta_H$, Tajima's
$D$, Fay \& Wu's $H$, the maximum DAF estimator, segregating-site
counts, and singleton counts; a two-population kernel additionally
produces Hudson's $F_{ST}$, Weir--Cockerham $F_{ST}$, $d_{XY}$, and
$d_a$. Statistics that depend on haplotype-pattern hashing (Garud's
$H_1$, $H_{12}$, $H_{123}$, $H_2/H_1$) and on scatter-add reductions
(mean nSL, the DAF histogram, $\mu_{\mathrm{SFS}}$, per-window $Z_{nS}$
and $\omega$, $\mu_{\mathrm{LD}}$, SNP-distance moments) reuse the same
window-binning infrastructure through dedicated kernel families.

\subsubsection*{Memory-aware chunking}

A single-launch fused kernel materializes a transposed working copy of
the matrix---an 8-bit integer or 32-bit float array, depending on the
kernel---that can exceed GPU memory for large datasets. \pggpu projects
this working-set size at runtime from a free-memory query and routes
the workload through a chunked variant of the same kernels whenever it
would exceed 70\,\% of free GPU memory. The chunked path tiles the
variant axis, calls the same kernels per chunk, and accumulates partial
sums to the final result, producing numerically identical output to the
single-pass path. This was sufficient to keep the full
10.9-million-variant Ag1000G arm on a single 80\,GB A100 with no user
tuning.

\subsubsection*{Missing data, accessibility, and span normalization}

Missing data and accessibility enter the single-locus ratios through
separate terms. Every diversity, divergence, and SFS estimator uses a
per-site count of non-missing samples in the numerator, so
allele-frequency estimates and their variance terms remain unbiased
when the called-sample fraction varies across sites. For the
denominator, an optional \texttt{AccessibleMask}---a boolean array over
genomic coordinates, supplied directly or loaded from a BED
file---replaces the nominal window width with the callable-bases count
in each window. Without such a mask, gaps, low-mappability regions, and
repeat-masked positions enter the denominator and silently deflate
per-base values for length-normalized statistics ($\pi$, $\theta_W$,
$d_{XY}$ per site) in poorly accessible regions, distorting windowed
scans even when every per-site number is itself unbiased. The per-site
valid-sample-count handling keeps the diversity estimator unbiased
across missing-data rates from 0 to 60\,\% (Supplementary
Figure~\ref{fig:missing_bias}). The convention matches
\scikitallel's windowed diversity and divergence functions, which
use the same per-site valid-count term in the numerator and accept
an accessibility array to mask sites that are flagged for exclusion.

\subsubsection*{Two-locus statistics}

Two-locus (linkage-disequilibrium) statistics are the most expensive
members of the catalog and the ones \pggpu accelerates most, since the
pairwise computation over variant sites is inherently
$O(n_{\mathrm{var}}^2)$. \pggpu factors every two-locus statistic
through a single count-based kernel: for each pair of sites it forms
the four haplotype-pair counts $(n_{AB}, n_{Ab}, n_{aB}, n_{ab})$ and
the per-pair non-missing sample count, and a kernel with one thread per
site pair emits any requested subset of $D^2$, $Dz$, $\pi_2$, $r^2$,
$D'$, and Pearson $r$ in double precision. Windowed two-locus summaries
(Kelly's $Z_{nS}$, Kim \& Nielsen's $\omega$, RAiSD's $\mu_{\mathrm{LD}}$)
reuse the window-binning machinery of the fused kernels; $Z_{nS}$
additionally offers a sampled-pairs estimator that trades the full
$O(n_{\mathrm{var}}^2)$ sum for a sub-quadratic profile
(Figure~\ref{fig:scaling}). The same count-based path underlies the
\moments-LD drop-in replacement (below), where the largest single
speedup ($\sim$1{,}750$\times$) is realized.

\subsubsection*{Interfacing with downstream tools}

\pggpu provides a drop-in replacement for the LD-parsing module of
\moments \citep{ragsdale2019}: the \texttt{compute\_ld\_statistics}
function consumes the same \vcf / recombination-map / population-file
inputs and returns a dictionary in \moments' canonical schema
\texttt{(bins, sums, stats, pops)}, so existing infrastructure
(\texttt{moments.\allowbreak LD.\allowbreak Parsing.\allowbreak bootstrap\_data}, \texttt{optimize\_log\_lbfgsb},
and uncertainty quantification) accepts the output unchanged. Local PCA
/ \lostruct \citep{li2019} is implemented as a single batched-Gram
eigendecomposition over windowed views of the matrix; the subsequent
inter-window MDS and corner-detection steps follow \citet{li2019}.

\clearpage

\subsection*{Extended feature description}

The \pggpu API mirrors the category organization of
Table~\ref{tab:features}: each category corresponds to a Python module
(\texttt{pg\_gpu.diversity}, \texttt{pg\_gpu.divergence},
\texttt{pg\_gpu.selection}, and so on), and each statistic is a
function that takes a \texttt{HaplotypeMatrix} or
\texttt{GenotypeMatrix} as its first argument and returns a \numpy
array. Population labels assigned once via \texttt{load\_pop\_file} are
accepted by every statistic that needs them, i.e.,
\texttt{divergence.fst\_hudson(h, "pop1", "pop2")} and
\texttt{diversity.tajimas\_d(h, population="pop1")} share the same
calling convention. Where a statistic is defined for both phased and
diploid input, the same call dispatches on container type:
\texttt{selection.garud\_h} computes haplotype-frequency $H$ statistics
on a \texttt{HaplotypeMatrix} and the unphased diploid analogue on a
\texttt{GenotypeMatrix} without separate method names. Missing data and
accessibility are handled automatically.

Beyond the classic estimators, \pggpu exposes the
Schrider et al.~\citep{schrider2018} distance-based two-population
statistics ($Z_x$, $dd$, $dd$ rank) and the RAiSD
$\mu_{\mathrm{LD}}$ component~\citep{alachiotis2018} as ordinary
functions; Patterson's $f$-statistics~\citep{patterson2012} return
per-site numerators and denominators separately so that ratio-of-sums
standard errors can be computed correctly with resampling. The
pairwise-distance moments of \citet{schrider2018} and the
diplotype-frequency spectrum used for sweep-classifier features are
also exposed directly.

The \texttt{FrequencySpectrum} class provides a direct interface to the
Achaz~\citep{achaz2009} weighted-SFS formulation, supporting arbitrary
$\theta$ estimators or neutrality tests outside the named catalog. A
\texttt{FrequencySpectrum} built from a \texttt{HaplotypeMatrix}
computes the SFS once on the GPU and then exposes three operations:
\texttt{theta(weights)} evaluates any linear estimator
$\hat\theta = \sum_\xi w_\xi \xi p_\xi$ given either a name from the
built-in dictionary (\texttt{pi}, \texttt{watterson}, \texttt{theta\_h},
\texttt{theta\_l}, \texttt{eta1}, etc.) or a user-supplied callable
that returns a weight vector; \texttt{neutrality\_test} pairs any two
estimators into a normalized $D$-style test statistic with a variance
computed from the cached $\alpha_n,\beta_n$ coefficients of
\citet{achaz2009}, so a custom neutrality test built around an unusual
weighting (e.g., higher-order powers of the frequency spectrum)
is one call rather than a hand-derived variance formula; and
\texttt{project} produces a hypergeometrically sub-sampled SFS at any
target sample size~\citep{gutenkunst2009}, allowing comparisons across
populations with different sample sizes and bridging to
demographic-inference pipelines that operate on a projected SFS.

Almost every statistic in Table~\ref{tab:features} is also accessible
through \texttt{windowed\_analysis}, a single entry point for genome
scans that takes a list of statistic names, a window definition (in
either base pairs or SNPs), and a list of populations, and returns a
\texttt{pandas.DataFrame}~\citep{mckinney2010} with one row per window.
Internally, \texttt{windowed\_analysis} routes the request through the
fused-kernel families described above, so a
typical scan that asks for $\pi$, $\theta_W$, Tajima's $D$, $F_{ST}$,
$d_{XY}$, Garud's $H_{12}$, and mean nSL---seven statistics spanning
four kernel families---runs in a small constant number of GPU launches
over the matrix rather than seven independent passes. Local PCA /
\lostruct~\citep{li2019} is exposed through the same interface and
returns a \texttt{LocalPCAResult} with per-window eigenvalues and
eigenvectors, which the companion functions \texttt{pc\_dist} and
\texttt{corners} transform into the inter-window distance matrix and
outlier-window identifications shown in
Figure~\ref{fig:lostruct_ag1000g}.

General-purpose \texttt{block\_jackknife} and \texttt{block\_bootstrap}
estimators~\citep{busing1999, efron1993} provide calibrated standard
errors and confidence intervals for any scalar aggregate of windowed
output. Both accept a tuple of (numerator, denominator) block arrays in
addition to a single array, so ratio-of-sums quantities such as
Patterson's $f_3$, $D$, and $F_{ST}$ are resampled with the same block
assignments on the two sides of the ratio without manual bookkeeping.

\clearpage

\setlength{\LTcapwidth}{\textwidth}
\small
\begin{longtable}{@{}>{\raggedright\arraybackslash}p{0.85in} >{\raggedright\arraybackslash}p{2.75in} >{\raggedright\arraybackslash}p{2.10in}@{}}
\caption{Complete catalog of statistics implemented in \pggpu.}
\label{tab:full_catalog} \\
\toprule
\textbf{Category} & \textbf{Statistic} & \textbf{Reference} \\
\midrule
\endfirsthead
\toprule
\textbf{Category} & \textbf{Statistic} & \textbf{Reference} \\
\midrule
\endhead
\midrule
\multicolumn{3}{r}{\textit{Continued on next page}} \\
\bottomrule
\endfoot
\bottomrule
\endlastfoot

% Diversity & neutrality tests
\multirow{15}{*}{Diversity} & Nucleotide diversity ($\pi$) & \citep{nei1979} \\
 & Watterson's $\theta_W$ & \citep{watterson1975} \\
 & Fay \& Wu's $\theta_H$ & \citep{fay2000} \\
 & Zeng's $\theta_L$ & \citep{zeng2006} \\
 & Tajima's $D$ & \citep{tajima1989} \\
 & Fay \& Wu's $H$, normalized $H^*$ & \citep{zeng2006} \\
 & Zeng's $E$ & \citep{zeng2006} \\
 & Zeng's $DH$ joint test & \citep{zeng2006} \\
 & Segregating sites, singletons & \citep{fu1995} \\
 & Haplotype diversity, count & \citep{nei1987} \\
 & Expected, observed heterozygosity & \citep{nei1973} \\
 & Wright's $F$ & \citep{wright1951} \\
 & Max DAF, DAF histogram & -- \\
 & RAiSD $\mu_{\mathrm{VAR}}$, $\mu_{\mathrm{SFS}}$ & \citep{alachiotis2018} \\
 & Diplotype frequency spectrum & -- \\
\midrule

% Divergence
\multirow{10}{*}{Divergence} & $F_{ST}$ (Hudson) & \citep{hudson1992} \\
 & $F_{ST}$ (Weir--Cockerham) & \citep{weir1984} \\
 & $F_{ST}$ (Nei) / $G_{ST}$ & \citep{nei1973} \\
 & $d_{XY}$ & \citep{nei1987} \\
 & $D_a$ (net divergence) & \citep{nei1979} \\
 & PBS & \citep{yi2010} \\
 & $S_{nn}$ & \citep{hudson2000} \\
 & $G_{\min}$ & \citep{geneva2015} \\
 & $d_d$, $d_d$ rank & \citep{schrider2018} \\
 & $Z_x$ & \citep{schrider2018} \\
\midrule

% Linkage disequilibrium
\multirow{8}{*}{LD} & $r$, $r^2$ & \citep{hill1968} \\
 & $D^2$, $D_z$, $\pi_2$ & \citep{ragsdale2020} \\
 & $\sigma_d^2$ (unbiased) & \citep{ragsdale2020} \\
 & Kelly's $Z_{nS}$ & \citep{kelly1997} \\
 & Kim \& Nielsen's $\omega$ & \citep{kim2004} \\
 & $\mu_{\mathrm{LD}}$ (RAiSD) & \citep{alachiotis2018} \\
 & Fused windowed LD & -- \\
 & \moments drop-in LD parsing & -- \\
\midrule

% Selection scans
\multirow{6}{*}{Selection} & iHS & \citep{voight2006} \\
 & nSL & \citep{ferrer-admetlla2014} \\
 & XP-EHH & \citep{sabeti2007} \\
 & XP-nSL & \citep{szpiech2021} \\
 & $H_1$, $H_{12}$, $H_{123}$, $H_2/H_1$ & \citep{garud2015} \\
 & EHH decay & \citep{sabeti2002} \\
\midrule

% SFS
\multirow{4}{*}{SFS} & Unfolded SFS, folded SFS & -- \\
 & Scaled SFS, folded-scaled SFS & -- \\
 & Joint (2-pop) SFS, folded joint SFS & -- \\
 & Scaled joint SFS, folded-scaled joint SFS & -- \\
\midrule

% Admixture / F-statistics
\multirow{5}{*}{Admixture} & Patterson's $F_2$ & \citep{patterson2012} \\
 & Patterson's $F_3$ & \citep{patterson2012} \\
 & Patterson's $D$ (ABBA-BABA) & \citep{patterson2012} \\
 & Windowed $F_3$, $D$ & \citep{patterson2012} \\
 & Block-jackknife $F_3$, $D$ standard error & \citep{busing1999} \\
\midrule

% Dimensionality reduction
\multirow{5}{*}{Dim.\ reduction} & PCA & \citep{patterson2006} \\
 & Randomized PCA & \citep{halko2011} \\
 & PCoA (classical MDS) & -- \\
 & \shortstack[l]{Local PCA / \lostruct \\ ~~(windowed eigendecomposition)} &
     \citep{li2019} \\
 & \shortstack[l]{Corner detection (Welzl \\ ~~minimum enclosing circle)} &
     \citep{li2019} \\
\midrule

% Resampling
\multirow{2}{*}{Resampling} & Block jackknife (incl.\ unequal blocks) & \citep{busing1999} \\
 & Block bootstrap & \citep{efron1993} \\
\midrule

% Generalized theta framework
\multirow{3}{*}{Generalized $\theta$} & Weighted-SFS $\theta$ estimator & \citep{achaz2009} \\
 & \shortstack[l]{Paired neutrality tests \\ ~~with cached variance} &
    \citep{achaz2009} \\
 & SFS projection (hypergeometric) & \citep{gutenkunst2009} \\
\midrule

% Relatedness
\multirow{2}{*}{Relatedness} & GRM & \citep{yang2011} \\
 & IBS & \citep{purcell2007} \\
\midrule

% Distance distribution
\multirow{4}{*}{Distance dist.} & Pairwise Hamming distance & -- \\
 & Variance of distances & \citep{schrider2018} \\
 & Skewness of distances & \citep{schrider2018} \\
 & Excess kurtosis of distances & \citep{schrider2018} \\

\end{longtable}

\end{document}
